% ============================================================
% Master file for the standalone RSA Cryptosystem module.
% Assembles Lessons 1-3 (RSA foundations, mathematical attacks,
% implementation/protocol attacks) into one combined "book" with
% a single table of contents and each lecture as a "Lesson N: ..."
% chapter -- same pattern as the parent course's full-course-notes.tex,
% just scoped to this module's three lessons.
%
% How this works:
%   - Each lectureN-rsa-*.tex file is still a fully independent,
%     standalone document (own title page, own table of contents)
%     -- it gets its content from lectureN-rsa-*-body.tex via \input.
%   - This file \input's the exact same body files directly, wrapped
%     in \chapter{Lesson N: ...} headings, under the `book` class.
%
% Usage:
%   pdflatex rsa-course-notes.tex   (run twice -- compile_all.sh
%   in this folder does this automatically)
% ============================================================
\documentclass[11pt,a4paper,oneside]{book}
% ============================================================
% Shared preamble for the RSA module's lecture files
% (lectureN-rsa-*.tex and rsa-course-notes.tex). \input this once
% per document.
%
% Defines: packages and the theorem-style environments used
% throughout the lecture bodies (theorem, lemma, corollary,
% proposition, definition, example, remark, exercise). Deliberately
% plain formatting -- no decorative call-out boxes -- so the
% lectures read as clean, standard mathematical exposition.
%
% Numbering: theorem-family environments are numbered "within
% section" ([section]), so they read as e.g. "Theorem 2.3" inside
% a standalone lectureN.tex (plain \section counter) and as e.g.
% "Theorem 2.2.3" inside rsa-course-notes.tex (book class nests
% \thesection under \thechapter automatically) -- no manual
% renumbering needed between the two contexts.
% ============================================================
\usepackage[margin=1in]{geometry}
\usepackage{amsmath,amssymb,amsthm}
\usepackage{algorithm}
\usepackage{algpseudocode}
\usepackage{enumitem}
\usepackage{hyperref}
\usepackage{booktabs}
\usepackage{xcolor}
\usepackage{titlesec}

% Single shared definition of the RSA module's series title. Every
% document in this project (each standalone lectureN-*.tex wrapper,
% and the combined rsa-course-notes.tex) uses \seriestitle inside its
% own \title{...} command instead of retyping the words -- so the
% title text itself is defined in exactly ONE place. Change it here
% once and every document picks it up on its next compile.
\newcommand{\seriestitle}{The RSA Cryptosystem: Foundations, Attacks, and Implementation}

% Lets \ref (not hardcoded numbers) resolve section/theorem labels that
% live in a DIFFERENT lecture's .tex file -- e.g. Lesson 2 citing a
% theorem number from Lesson 1. Each lecture wrapper file declares which
% other lecture(s) it needs via \externaldocument{<other-lecture-basename>}
% right after % ============================================================
% Shared preamble for the RSA module's lecture files
% (lectureN-rsa-*.tex and rsa-course-notes.tex). \input this once
% per document.
%
% Defines: packages and the theorem-style environments used
% throughout the lecture bodies (theorem, lemma, corollary,
% proposition, definition, example, remark, exercise). Deliberately
% plain formatting -- no decorative call-out boxes -- so the
% lectures read as clean, standard mathematical exposition.
%
% Numbering: theorem-family environments are numbered "within
% section" ([section]), so they read as e.g. "Theorem 2.3" inside
% a standalone lectureN.tex (plain \section counter) and as e.g.
% "Theorem 2.2.3" inside rsa-course-notes.tex (book class nests
% \thesection under \thechapter automatically) -- no manual
% renumbering needed between the two contexts.
% ============================================================
\usepackage[margin=1in]{geometry}
\usepackage{amsmath,amssymb,amsthm}
\usepackage{algorithm}
\usepackage{algpseudocode}
\usepackage{enumitem}
\usepackage{hyperref}
\usepackage{booktabs}
\usepackage{xcolor}
\usepackage{titlesec}

% Single shared definition of the RSA module's series title. Every
% document in this project (each standalone lectureN-*.tex wrapper,
% and the combined rsa-course-notes.tex) uses \seriestitle inside its
% own \title{...} command instead of retyping the words -- so the
% title text itself is defined in exactly ONE place. Change it here
% once and every document picks it up on its next compile.
\newcommand{\seriestitle}{The RSA Cryptosystem: Foundations, Attacks, and Implementation}

% Lets \ref (not hardcoded numbers) resolve section/theorem labels that
% live in a DIFFERENT lecture's .tex file -- e.g. Lesson 2 citing a
% theorem number from Lesson 1. Each lecture wrapper file declares which
% other lecture(s) it needs via \externaldocument{<other-lecture-basename>}
% right after % ============================================================
% Shared preamble for the RSA module's lecture files
% (lectureN-rsa-*.tex and rsa-course-notes.tex). \input this once
% per document.
%
% Defines: packages and the theorem-style environments used
% throughout the lecture bodies (theorem, lemma, corollary,
% proposition, definition, example, remark, exercise). Deliberately
% plain formatting -- no decorative call-out boxes -- so the
% lectures read as clean, standard mathematical exposition.
%
% Numbering: theorem-family environments are numbered "within
% section" ([section]), so they read as e.g. "Theorem 2.3" inside
% a standalone lectureN.tex (plain \section counter) and as e.g.
% "Theorem 2.2.3" inside rsa-course-notes.tex (book class nests
% \thesection under \thechapter automatically) -- no manual
% renumbering needed between the two contexts.
% ============================================================
\usepackage[margin=1in]{geometry}
\usepackage{amsmath,amssymb,amsthm}
\usepackage{algorithm}
\usepackage{algpseudocode}
\usepackage{enumitem}
\usepackage{hyperref}
\usepackage{booktabs}
\usepackage{xcolor}
\usepackage{titlesec}

% Single shared definition of the RSA module's series title. Every
% document in this project (each standalone lectureN-*.tex wrapper,
% and the combined rsa-course-notes.tex) uses \seriestitle inside its
% own \title{...} command instead of retyping the words -- so the
% title text itself is defined in exactly ONE place. Change it here
% once and every document picks it up on its next compile.
\newcommand{\seriestitle}{The RSA Cryptosystem: Foundations, Attacks, and Implementation}

% Lets \ref (not hardcoded numbers) resolve section/theorem labels that
% live in a DIFFERENT lecture's .tex file -- e.g. Lesson 2 citing a
% theorem number from Lesson 1. Each lecture wrapper file declares which
% other lecture(s) it needs via \externaldocument{<other-lecture-basename>}
% right after \input{lecture-preamble}. This requires the referenced
% lecture to have been compiled at least once already (so its .aux
% exists) -- compile_all.sh handles this automatically by building the
% lectures in order (1, then 2, then 3). If you compile a single later
% lecture in isolation before ever compiling the earlier one(s) it cites,
% those specific cross-references will show as "??" until you compile
% the earlier lecture(s) once.
\usepackage{xr-hyper}

\newtheorem{theorem}{Theorem}[section]
\newtheorem{lemma}[theorem]{Lemma}
\newtheorem{corollary}[theorem]{Corollary}
\newtheorem{proposition}[theorem]{Proposition}
\newtheorem{definition}[theorem]{Definition}
\newtheorem{example}[theorem]{Example}
\newtheorem{remark}[theorem]{Remark}
\newtheorem{exercise}[theorem]{Exercise}

\hypersetup{
    colorlinks=true,
    linkcolor=blue!50!black,
    citecolor=blue!50!black,
    urlcolor=blue!50!black,
    pdfpagemode=UseNone
}
. This requires the referenced
% lecture to have been compiled at least once already (so its .aux
% exists) -- compile_all.sh handles this automatically by building the
% lectures in order (1, then 2, then 3). If you compile a single later
% lecture in isolation before ever compiling the earlier one(s) it cites,
% those specific cross-references will show as "??" until you compile
% the earlier lecture(s) once.
\usepackage{xr-hyper}

\newtheorem{theorem}{Theorem}[section]
\newtheorem{lemma}[theorem]{Lemma}
\newtheorem{corollary}[theorem]{Corollary}
\newtheorem{proposition}[theorem]{Proposition}
\newtheorem{definition}[theorem]{Definition}
\newtheorem{example}[theorem]{Example}
\newtheorem{remark}[theorem]{Remark}
\newtheorem{exercise}[theorem]{Exercise}

\hypersetup{
    colorlinks=true,
    linkcolor=blue!50!black,
    citecolor=blue!50!black,
    urlcolor=blue!50!black,
    pdfpagemode=UseNone
}
. This requires the referenced
% lecture to have been compiled at least once already (so its .aux
% exists) -- compile_all.sh handles this automatically by building the
% lectures in order (1, then 2, then 3). If you compile a single later
% lecture in isolation before ever compiling the earlier one(s) it cites,
% those specific cross-references will show as "??" until you compile
% the earlier lecture(s) once.
\usepackage{xr-hyper}

\newtheorem{theorem}{Theorem}[section]
\newtheorem{lemma}[theorem]{Lemma}
\newtheorem{corollary}[theorem]{Corollary}
\newtheorem{proposition}[theorem]{Proposition}
\newtheorem{definition}[theorem]{Definition}
\newtheorem{example}[theorem]{Example}
\newtheorem{remark}[theorem]{Remark}
\newtheorem{exercise}[theorem]{Exercise}

\hypersetup{
    colorlinks=true,
    linkcolor=blue!50!black,
    citecolor=blue!50!black,
    urlcolor=blue!50!black,
    pdfpagemode=UseNone
}


% Same header customization as the parent course's master file:
% chapter headings read as plain "Lesson N: Title", no automatic
% "Chapter N" prefix, and a plain page style (long titles would
% overflow the book class's default running header).
\titleformat{\chapter}[block]{\normalfont\huge\bfseries}{}{0pt}{}
\titlespacing*{\chapter}{0pt}{-20pt}{30pt}
\pagestyle{plain}

\title{\textbf{Applied Information Security}\\[6pt]
\large The RSA Cryptosystem: Foundations, Attacks, and Implementation}
\author{Dr.\ Faisal Aslam}
\date{}

\begin{document}
\maketitle
\tableofcontents

\chapter{Lesson 1: The RSA Cryptosystem I --- Foundations, Correctness, and Factoring}
\begin{center}
\large\itshape Euler's Theorem, the Chinese Remainder Theorem, Key Generation, and Classical Factoring Algorithms
\end{center}
\vspace{0.4cm}
\input{lecture1-rsa-foundations-body}

\chapter{Lesson 2: The RSA Cryptosystem II --- Mathematical Attacks}
\begin{center}
\large\itshape Wiener's Attack, Håstad's Broadcast Attack, Common Modulus Attacks, Franklin--Reiter, and Coppersmith's Method
\end{center}
\vspace{0.4cm}
\vspace{0.5cm}
\subsection*{Roadmap}
\begin{enumerate}[label=(\arabic*)]
\item Wiener's Attack: recovering a small private exponent $d$ from $(n,e)$ alone, via continued fractions
\item Håstad's Broadcast Attack: recovering $m$ from low-exponent encryptions to multiple recipients, via CRT
\item Common Modulus Attack: breaking RSA when two parties share $n$ but use different, coprime exponents
\item Cycling Attack: a conceptual (impractical) attack via the multiplicative order of $c$
\item The Franklin--Reiter Related-Message Attack: breaking two related plaintexts under the same low exponent, via polynomial GCD
\item Coppersmith's Method: an overview of small-root finding via lattice reduction (the LLL algorithm)
\end{enumerate}
\textbf{What every attack in this lecture has in common:} none of them factor $n$. Each exploits a specific way RSA was \emph{used} — a too-small $d$, a too-small $e$ combined with related or repeated messages, exponent/modulus reuse, or the sheer algebraic structure of modular exponentiation. This is the point made in Lesson 1: breaking RSA is not the same as solving the RSA problem in general, and it is very much not the same as factoring.

\newpage
\section{Wiener's Attack: Small Private Exponent}

\subsection{Motivation}
CRT-RSA decryption (Lesson 1, Section~\ref{sec:crt-speedup}) is faster with smaller $d$. This creates a temptation: why not simply choose a small $d$ to speed up decryption, and derive $e = d^{-1} \bmod \varphi(n)$ afterward? Wiener showed in 1990 that this is catastrophic: if $d$ is small enough, the \emph{entire private key} is recoverable from the public key alone, in a fraction of a second, using nothing more exotic than the Euclidean algorithm.

\medskip\noindent\textbf{The key insight, in one sentence, before any machinery.} The public key $(n,e)$ is, on its face, just two large numbers with no visible structure. But if $d$ is small, it turns out that the fraction $e/n$ — which the attacker already knows exactly — is secretly an extremely good rational approximation of a related fraction $k/d$ involving the private $d$. ``Extremely good approximations of a fraction with a small denominator'' are rare and easy to enumerate, so the attacker can simply search that short list rather than guessing $d$ directly. Continued fractions are the tool that makes ``search the short list of extremely good approximations'' precise and efficient; we build that tool first, then use it to attack RSA.

\subsection{Continued Fractions Refresher}

\medskip\noindent\textbf{What a continued fraction is, and why anyone invented it.} Every rational number can be written as a whole part plus a fraction, and that fraction can itself be flipped and split the same way, recursively:
\[
\frac{30}{7} = 4 + \frac{2}{7} = 4 + \cfrac{1}{7/2} = 4 + \cfrac{1}{3 + \tfrac12} = 4 + \cfrac{1}{3 + \cfrac{1}{2}}.
\]
Stopping this process at each stage and simplifying back into a single fraction gives a sequence of increasingly accurate rational approximations to the original number — $4$, then $4+\frac13 = \frac{13}{3}$, then the exact value $\frac{30}{7}$ itself. This is the entire idea; the definitions below just make it precise and give a fast way to compute it (the same Euclidean algorithm you already know from computing $\gcd$s and modular inverses).

\begin{definition}[Continued fraction expansion]
Every rational number $\frac{x}{y}$ has a unique finite continued fraction expansion
\[
\frac{x}{y} = a_0 + \cfrac{1}{a_1 + \cfrac{1}{a_2 + \cfrac{1}{\ddots + \cfrac{1}{a_k}}}} =: [a_0; a_1, a_2, \dots, a_k],
\]
computed by the Euclidean algorithm: $a_0 = \lfloor x/y \rfloor$, then recurse on $y / (x - a_0 y)$.
\end{definition}

\begin{example}[Continued fraction expansion, small warm-up]
Expand $\frac{30}{7}$. Step 1: $a_0 = \lfloor 30/7\rfloor = 4$, remainder $30 - 4\cdot7 = 2$; recurse on $7/2$. Step 2: $a_1 = \lfloor 7/2 \rfloor = 3$, remainder $7-3\cdot2=1$; recurse on $2/1$. Step 3: $a_2 = \lfloor 2/1 \rfloor = 2$, remainder $0$ — stop. So $\frac{30}{7} = [4;3,2]$, matching the hand-expansion above. Notice this is \emph{exactly} the Euclidean algorithm run on $(30,7)$, just recording each quotient instead of only the final $\gcd$.
\end{example}

\begin{definition}[Convergents]
The $i$-th \textbf{convergent} $\frac{h_i}{k_i}$ of $[a_0;a_1,\dots]$ is the value of the truncated expansion $[a_0;a_1,\dots,a_i]$, computed by the recurrence
\[
h_i = a_i h_{i-1} + h_{i-2}, \qquad k_i = a_i k_{i-1} + k_{i-2}, \qquad (h_{-1}{=}1, h_{-2}{=}0,\ k_{-1}{=}0, k_{-2}{=}1).
\]
\end{definition}

\begin{example}[Convergents of the warm-up example]
Continuing $\frac{30}{7}=[4;3,2]$: the convergents are
\[
\frac{h_0}{k_0}=\frac{4}{1}=4, \qquad \frac{h_1}{k_1}=\frac{3\cdot4+1}{3\cdot1+0}=\frac{13}{3}\approx4.333, \qquad \frac{h_2}{k_2}=\frac{2\cdot13+4}{2\cdot3+1}=\frac{30}{7}\approx4.286.
\]
Each convergent is a \emph{better} rational approximation to $\frac{30}{7}=4.2857\dots$ than the last, and the final one is exact. This ``successively better approximation'' behavior, not just the recursive bookkeeping, is what Wiener's attack actually exploits below.
\end{example}

\begin{theorem}[Legendre's theorem on best rational approximations]
If $\left|\frac{x}{y} - \frac{h}{k}\right| < \frac{1}{2k^2}$, then $\frac{h}{k}$ is a convergent of the continued fraction expansion of $\frac{x}{y}$.
\end{theorem}

\medskip\noindent\textbf{Why this matters for RSA, unpacked.} Read the theorem as a converse: it does \emph{not} say every convergent is a great approximation (that direction is easy and unsurprising); it says every \emph{sufficiently good} approximation \emph{must already be} one of the (few!) convergents. There are astronomically many fractions $h/k$ with $k$ in some range, but only $O(\log y)$ convergents — one per step of the Euclidean algorithm on $(x,y)$. So instead of searching an exponentially large space of candidate fractions, the theorem lets an attacker search a logarithmically short, explicitly computable list, \emph{provided} they can first show that the fraction they are looking for happens to approximate a known quantity unusually well. That is exactly the situation Wiener's theorem engineers next.

\subsection{Wiener's Theorem}

\medskip\noindent\textbf{The plan, stated before the algebra.} The attacker knows $(n,e)$ and wants $d$. The relationship defining $d$ is $ed \equiv 1 \pmod{\varphi(n)}$, which involves the unknown $\varphi(n)$ as well as the unknown $d$ — two unknowns, seemingly hopeless. The trick is to show that this relationship, rewritten as an inequality about \emph{fractions}, forces $e/n$ (fully known) to sit unusually close to $k/d$ (fully unknown, but with small denominator $d$) — close enough that Legendre's theorem applies. Once that is shown, $k/d$ is guaranteed to appear among the convergents of $e/n$, which the attacker can simply compute and test one by one.

\begin{theorem}[Wiener, 1990]
Let $n = pq$ with $q < p < 2q$, and let $d < \frac{1}{3} n^{1/4}$. Then given only the public key $(n,e)$ with $ed \equiv 1 \pmod{\varphi(n)}$, $d$ can be recovered in polynomial time.
\end{theorem}

\begin{proof}[Proof, step by step]
\textbf{Step 1: turn the congruence into an equation, and identify the target fraction.} Since $ed \equiv 1 \pmod{\varphi(n)}$, there is a positive integer $k$ with
\[
ed = 1 + k\varphi(n).
\]
(One can check $k<d$: since $e<\varphi(n)$, we have $k = \frac{ed-1}{\varphi(n)} < \frac{ed}{\varphi(n)} < d$.) The quantity we will show is well-approximated by the known $e/n$ is precisely $k/d$ — a fraction built entirely from unknowns, but with denominator $d$, which is small by hypothesis.

\textbf{Step 2: measure how far apart $e/n$ and $k/d$ actually are.} Compute the difference directly:
\[
\left|\frac{e}{n} - \frac{k}{d}\right| = \left|\frac{ed - kn}{nd}\right| = \left|\frac{(1+k\varphi(n)) - kn}{nd}\right| = \frac{k\big(n - \varphi(n)\big) - 1}{nd}.
\]
The first equality is just putting both fractions over the common denominator $nd$; the second substitutes $ed = 1+k\varphi(n)$ from Step 1; the third regroups the terms. Everything now hinges on how large $n-\varphi(n)$ is.

\textbf{Step 3: bound $n - \varphi(n)$ using the balanced-primes assumption.} Recall $\varphi(n) = (p-1)(q-1) = n - (p+q) + 1$, so $n - \varphi(n) = p+q-1$. Using $q<p<2q$ (i.e.\ the two primes are within a factor of $2$ of each other — the normal case for RSA key generation) gives $p+q < 3q < 3\sqrt{n}\cdot\sqrt{2}$, and a slightly more careful accounting (using $q<\sqrt{n}$ and $p<2q$ together) gives the clean bound $n-\varphi(n) < 3\sqrt n$.

\textbf{Step 4: combine Steps 2 and 3, and invoke the hypothesis on $d$.} Substituting the Step 3 bound and $k<d$ into the Step 2 expression gives
\[
\left|\frac{e}{n}-\frac{k}{d}\right| = \frac{k(n-\varphi(n))-1}{nd} < \frac{k\cdot3\sqrt n}{nd} < \frac{3\sqrt n}{n} = \frac{3}{\sqrt n}.
\]
A short calculation (using $d<\frac13n^{1/4}$ to control the denominator $d^2$ against $\sqrt n$) tightens this into exactly the bound Legendre's theorem needs: $\left|\frac{e}{n}-\frac{k}{d}\right| < \frac{1}{2d^2}$.

\textbf{Step 5: invoke Legendre's theorem.} Since $\left|\frac{e}{n}-\frac{k}{d}\right|<\frac{1}{2d^2}$, Legendre's theorem guarantees $\frac{k}{d}$ is \emph{some} convergent of the continued fraction expansion of $\frac{e}{n}$ — a quantity the attacker already knows in full, from the public key alone, and whose convergents can be listed in $O(\log n)$ steps via the Euclidean algorithm.
\end{proof}

\medskip\noindent\textbf{What just happened, in one sentence.} A relationship between two \emph{unknowns} ($e$ and $d$'s hidden partner $k$) got converted, purely by algebra, into a statement that two \emph{known} numbers ($e$ and $n$) sit unusually close together as a fraction — and ``unusually close'' is exactly the trigger condition for Legendre's theorem to hand the attacker a short, explicit candidate list.

\begin{algorithm}[h]
\caption{Wiener's Attack}
\begin{algorithmic}[1]
\State Compute the continued fraction expansion of $e/n$
\State Compute all convergents $h_i/k_i$
\For{each convergent $k/d_{\text{cand}}$}
    \If{$k \ne 0$ and $(e \cdot d_{\text{cand}} - 1) \bmod k = 0$}
        \State $\varphi_{\text{cand}} \gets (e \cdot d_{\text{cand}} - 1)/k$
        \State Solve $x^2 - (n - \varphi_{\text{cand}} + 1)x + n = 0$ for integer roots $p, q$
        \If{integer roots exist and $pq = n$} \State \Return $d = d_{\text{cand}}$ (attack succeeded) \EndIf
    \EndIf
\EndFor
\end{algorithmic}
\end{algorithm}

\medskip\noindent\textbf{Reading the algorithm against the proof.} Line 3's convergents are exactly the candidates Step 5 guarantees contains the true $k/d$. Line 4 checks whether a candidate is \emph{consistent} with being the true $(k,d)$: if it is, then $\varphi_{\text{cand}} = (ed_{\text{cand}}-1)/k$ recomputes what $\varphi(n)$ would have to be. Line 6 then uses $\varphi_{\text{cand}}$ and $n$ together to solve for $p,q$ directly — since $p+q = n-\varphi_{\text{cand}}+1$ and $pq=n$ are two equations in two unknowns, giving a quadratic. If that quadratic has integer roots multiplying to $n$, the candidate was correct; if not, move to the next convergent. Because only $O(\log n)$ convergents exist, this whole search is fast even though it looks like a brute-force loop.

\begin{example}[Wiener's attack, fully worked]
Let $p = 1223$, $q = 1987$, so $n = 2{,}430{,}101$ and $\varphi(n) = 2{,}426{,}892$. Suppose the key-generator (unwisely) chose a small private exponent $d = 25$ for decryption speed. Then $e = d^{-1} \bmod \varphi(n) = 2{,}135{,}665$ — this is the only thing published.

\textbf{Check the bound}: $\frac{1}{3}n^{1/4} \approx 13.16$. Wait — $d=25 > 13.16$, so strictly this instance sits outside Wiener's original (conservative) bound; we use it anyway because the attack succeeds in practice on this instance regardless (the bound is a sufficient, not tight necessary, condition — later refinements by Boneh–Durfee push the provable bound up to $d < n^{0.292}$, which comfortably covers this example). This is worth flagging explicitly to students: textbook bounds are usually conservative sufficient conditions, not sharp thresholds.

\textbf{Run the attack}: the continued fraction expansion of $e/n = 2135665/2430101$ is
\[
[0; 1, 7, 3, 1, 17, 1, 1, 2, 1, 2, 211].
\]
Its convergents $h_i/k_i$ include (among others) $\frac{22}{25}$. Testing $k=22, d_{\text{cand}}=25$: $(e \cdot 25 - 1) \bmod 22 = 0$, giving $\varphi_{\text{cand}} = 2{,}426{,}892$ — exactly right. Solving the resulting quadratic for $p,q$ recovers $p=1223, q=1987$ \textbf{exactly}. The attacker recovered the \emph{entire private key} — not just $d$, but the factorization of $n$ itself — using only the public key $(n,e)$ and a handful of Euclidean-algorithm steps.
\end{example}

\medskip\noindent\textbf{Practical consequence.} \textbf{Never} choose $d$ for efficiency and derive $e$ afterward. Standard practice is the reverse: fix a large, well-understood $e$ (e.g., $65537$) and compute $d = e^{-1} \bmod \varphi(n)$, which is essentially always large. If CRT-RSA speed is the actual goal, achieve it via the CRT decomposition of Lesson 1 (Section~\ref{sec:crt-speedup}) instead — that optimization does not shrink $d$ itself, only the size of the exponents used \emph{internally} during decryption.

\newpage
\section{Håstad's Broadcast Attack}

\subsection{Setup}
Suppose a sender uses a \textbf{small public exponent} (e.g., $e=3$) and sends the \emph{same} message $m$, encrypted under $e$, to $e$ or more different recipients with pairwise coprime moduli $n_1, \dots, n_e$ (their own independent RSA keys):
\[
c_i = m^e \bmod n_i, \qquad i = 1, \dots, e.
\]

\medskip\noindent\textbf{Why this should feel dangerous, before the proof.} Recall the ``no modular wraparound'' example from Lesson 1: if $m^e < n$, then $c=m^e$ is simply the integer $m^e$ with no reduction ever taking place, and an ordinary integer root recovers $m$ instantly. A single well-generated RSA key avoids this by making $n$ enormous compared to any realistic $m^e$. Håstad's insight is that broadcasting the \emph{same} message to \emph{several} recipients recreates exactly this vulnerability one level up: even though each individual $c_i = m^e \bmod n_i$ wraps around plenty, the Chinese Remainder Theorem lets an eavesdropper combine all $e$ ciphertexts into a single value modulo the much larger product $N=n_1\cdots n_e$ — and against that enormous combined modulus, $m^e$ itself might not wrap around at all. Three small worlds, stitched together by CRT, behave like one giant world in which the ``no wraparound'' flaw resurfaces.

\begin{theorem}[Håstad's broadcast attack, $e=3$ identical-message case]
If $m < \min(n_1,n_2,n_3)$, an eavesdropper who observes $c_1,c_2,c_3$ (and the public moduli) recovers $m$ in polynomial time, without knowing any private key.
\end{theorem}

\begin{proof}[Proof, step by step]
\textbf{Step 1: combine the three congruences with CRT.} The eavesdropper knows $c_1,c_2,c_3$ and the public $n_1,n_2,n_3$ (pairwise coprime), and wants a single $x$ with $x\equiv c_i \pmod{n_i}$ for $i=1,2,3$. By CRT (Lesson 1, Theorem~\ref{thm:crt}, extended to three coprime moduli the same way it works for two), this system has a \emph{unique} solution $C$ modulo $N=n_1n_2n_3$, and — crucially — CRT is constructive: $C$ is computable directly from the public $c_i,n_i$ values using the reconstruction formula from Lesson 1, no private key involved anywhere.

\textbf{Step 2: notice $m^3$ itself satisfies the same three congruences.} Since $c_i \equiv m^3 \pmod{n_i}$ for each $i$ by definition of encryption, the value $m^3$ is \emph{also} a solution to the system in Step 1. By the uniqueness half of CRT, there is only one solution modulo $N$ — so $C \equiv m^3 \pmod N$.

\textbf{Step 3: upgrade ``congruent mod $N$'' to ``equal as integers''.} This is the step where the size assumption $m<\min(n_1,n_2,n_3)$ earns its keep: it gives $m^3 < n_1n_2n_3 = N$. Since $C$ (as CRT constructs it) is itself a specific representative in $\{0,\dots,N-1\}$, and $m^3$ is a nonnegative integer strictly less than $N$ that is congruent to $C$ modulo $N$, the two must be \emph{exactly} the same integer, not merely congruent — there is no room for them to differ by a multiple of $N$. So $C = m^3$ holds as an honest equation over $\mathbb{Z}$, not just modulo $N$.

\textbf{Step 4: take an ordinary integer cube root.} Since $C$ is now known to be exactly $m^3$ (an integer, not a residue), an ordinary integer cube-root computation — no modular arithmetic, no private key — recovers $m$.
\end{proof}

\medskip\noindent\textbf{The one-sentence summary.} Three separate ``small worlds'' (mod $n_1$, mod $n_2$, mod $n_3$), each individually opaque, get stitched by CRT into one ``big world'' (mod $N=n_1n_2n_3$) in which the encryption exponent's usual protection — reduction mod a modulus much larger than the message — silently stops applying.

\begin{example}[Håstad's attack, $e=3$, fully worked]
Let three (toy) recipients have moduli $n_1 = 143$, $n_2 = 323$, $n_3 = 667$ (pairwise coprime by construction: $143=11\cdot13$, $323=17\cdot19$, $667=23\cdot29$ — no shared prime factors). A sender broadcasts $m=17$ to all three under $e=3$:
\[
c_1 = 17^3 \bmod 143 = 51, \qquad c_2 = 17^3 \bmod 323 = 68, \qquad c_3 = 17^3 \bmod 667 = 244.
\]
An eavesdropper computes $N = n_1n_2n_3 = 30{,}808{,}063$ and solves the CRT system to get $C = 4913$. Since $m^3 = 17^3 = 4913 < N$, this $C$ is the true integer value of $m^3$ (no wraparound occurred). Taking the integer cube root: $\sqrt[3]{4913} = 17 = m$. \textbf{No private key was ever touched.}
\end{example}

\begin{remark}[Generalizing beyond identical messages]
The identical-message case above is the simplest illustration. If instead each recipient gets a \emph{linearly related} message (e.g., $m_i = a_i m + b_i$ for known public $a_i,b_i$ — a common real-world pattern, since padding schemes or per-recipient headers often relate messages this way), the same underlying idea extends via a more general CRT-plus-resultant construction. This generalization is due to Håstad, and its \emph{unknown-relation} extension (where even the relationship between messages is not fully known, only that it is small) is precisely Coppersmith's short-pad attack, previewed at the end of this lecture and building on the lattice technique in Section~\ref{sec:coppersmith}.
\end{remark}

\medskip\noindent\textbf{Practical consequence.} \textbf{Never} use a small public exponent to send the same (or simply related) message to multiple recipients without randomized padding. OAEP (Lesson 3) defeats this attack by making every ciphertext of the same message under the same key look independently random, so no algebraic relationship exists across ciphertexts for CRT to exploit.

\newpage
\section{Common Modulus Attack}

\subsection{Setup}
A tempting (and wrong) efficiency idea: have a central authority generate a single modulus $n$, and give each user their own $(e_i, d_i)$ pair on the \emph{same} $n$, to save on key-generation cost. Suppose a message $m$ is (perhaps accidentally, perhaps because the sender does not know which key to use) encrypted under two different users' exponents $e_1, e_2$ on the shared $n$, with $\gcd(e_1,e_2)=1$:
\[
c_1 = m^{e_1} \bmod n, \qquad c_2 = m^{e_2} \bmod n.
\]

\medskip\noindent\textbf{Why this should feel dangerous, before the proof.} Think about what $\gcd(e_1,e_2)=1$ actually buys you: by the Extended Euclidean Algorithm (the same tool from Lesson 1 that computes $d=e^{-1}\bmod\varphi(n)$), there exist integers $a,b$ with $ae_1+be_2=1$. That is a way of writing the number $1$ itself as a combination of $e_1$ and $e_2$. Since $c_1$ and $c_2$ are both just $m$ raised to a power, and exponents combine by \emph{multiplying} when you raise a power to a power, ``combining $e_1$ and $e_2$ to make $1$'' translates directly into ``combining $c_1$ and $c_2$ to make $m^1=m$'' — no private key needed at any point, because everything used here ($e_1,e_2,c_1,c_2,n$) is public.

\begin{theorem}[Common modulus attack]
If $\gcd(e_1,e_2)=1$, an eavesdropper who knows only $(n, e_1, e_2, c_1, c_2)$ — no private key at all — recovers $m$.
\end{theorem}

\begin{proof}[Proof, step by step]
\textbf{Step 1: find Bézout coefficients for $e_1,e_2$.} Since $\gcd(e_1,e_2)=1$, the Extended Euclidean Algorithm produces integers $a,b$ (exactly one of which will typically be negative) with
\[
ae_1+be_2=1.
\]

\textbf{Step 2: raise each ciphertext to the matching coefficient, and multiply.} Compute $c_1^a c_2^b \bmod n$ and expand using the definitions $c_1=m^{e_1}$, $c_2=m^{e_2}$:
\[
c_1^{a}c_2^{b} \equiv (m^{e_1})^a(m^{e_2})^b \equiv m^{ae_1}\cdot m^{be_2} \equiv m^{ae_1+be_2} \pmod n.
\]
The first equality substitutes the definitions of $c_1,c_2$; the second and third just apply the exponent laws $x^{mn}=(x^m)^n$ and $x^mx^n=x^{m+n}$ — nothing here uses any secret.

\textbf{Step 3: substitute Step 1's identity into the exponent.} Since $ae_1+be_2=1$ exactly (an integer identity, not just a congruence), the exponent in Step 2 collapses:
\[
m^{ae_1+be_2} = m^1 = m.
\]
Chaining Steps 2 and 3 gives $c_1^ac_2^b \equiv m \pmod n$ — the plaintext, recovered from two public ciphertexts and two public exponents alone.

\textbf{Step 4: handle a negative exponent.} If (say) $b<0$, then $c_2^b$ is not directly computable by repeated squaring in the usual sense; instead compute $c_2^{-1}\bmod n$ first via the Extended Euclidean Algorithm (this requires $\gcd(c_2,n)=1$, true with overwhelming probability for a random message — the exceptional case, $\gcd(c_2,n)\ne1$, would itself immediately leak a factor of $n$ via a single gcd computation, an even worse outcome for the key owner), and then raise that inverse to the power $|b|$ instead.
\end{proof}

\medskip\noindent\textbf{The one-sentence summary.} Two exponentiations of the \emph{same} secret message, under two exponents that are coprime, can always be algebraically recombined into the message itself, because coprimality is exactly the condition under which $1$ can be written as an integer combination of the two exponents.

\begin{example}[Common modulus attack, fully worked]
Let $n = 3233$ (our familiar $p=61,q=53$ toy modulus). Two users share $n$ with $e_1 = 17$, $e_2=7$ (coprime). A message $m=42$ is encrypted under both: $c_1 = 42^{17} \bmod 3233 = 2557$, $c_2 = 42^{7} \bmod 3233 = 240$.

Extended Euclid on $(17,7)$ gives $a=-2, b=5$ (check: $-2\cdot17 + 5\cdot7 = -34+35=1$). Since $a<0$, compute $c_1^{-1} \bmod n = 2160$ (via modular inverse), then
\[
m = (c_1^{-1})^{2} \cdot c_2^{5} \bmod n = 42.
\]
Recovered exactly — \textbf{using no private key whatsoever}, just the two public exponents and two ciphertexts.
\end{example}

\medskip\noindent\textbf{Practical consequence.} \textbf{Never} share a modulus $n$ across users or purposes, even with different exponents. This design pattern occasionally resurfaces in embedded/legacy systems trying to save key-generation cost, but it is fundamentally broken the moment any two exponents used on the shared modulus are coprime — which is the common case, since exponents are typically chosen small and prime.

\newpage
\section{Cycling Attack (Conceptual)}

\subsection{Setup}
This attack requires no assumption about $d$, $e$, or shared moduli — only encryption itself, and is presented mainly for conceptual completeness (it is essentially never practical at real key sizes, but it is instructive about the group-theoretic structure of RSA).

\medskip\noindent\textbf{Why this should feel familiar, before the proof.} Recall from Lesson 1 that repeatedly multiplying a group element $a$ by itself eventually returns to $1$ — that was the whole content of Euler's theorem, proved there by showing multiplication-by-$a$ just permutes the group. \emph{Repeated encryption} is exactly repeated exponentiation by the same fixed exponent $e$, so the same cyclic idea applies one level up: instead of asking ``how many multiplications by $a$ get back to $1$?'' we ask ``how many \emph{encryptions} (exponentiations by $e$) get back to the original ciphertext?'' Because encryption is invertible (that is the entire point of RSA), this cycle must eventually close, and — remarkably — the plaintext sits in plain view one step before it closes.

\begin{theorem}[Cycling attack]
Given only $(n,e,c)$ where $c = m^e \bmod n$, repeatedly compute $c_1 = c^e \bmod n$, $c_2 = c_1^{e} \bmod n$, and so on, until some $c_k = c$. Then $m = c_{k-1}$.
\end{theorem}

\begin{proof}[Proof, step by step]
\textbf{Step 1: express every term of the sequence as a single power of $m$.} Encrypting repeatedly under the public exponent is the same as encrypting-then-encrypting, so writing $c_0=c$, an easy induction gives $c_i \equiv m^{e^{i+1}} \pmod n$ for every $i\ge0$ (each step multiplies the current exponent of $m$ by another factor of $e$).

\textbf{Step 2: define $k$ as the point where the sequence first returns to its start.} Let $k$ be the smallest positive integer with $c_k \equiv c_0 \pmod n$. By Step 1, this says $m^{e^{k+1}} \equiv m^{e} \pmod n$.

\textbf{Step 3: cancel one factor of $e$ from the exponent, using invertibility.} Assuming $\gcd(m,n)=1$, the congruence $m^{e^{k+1}}\equiv m^e\pmod n$ is really a statement about the \emph{order} of $m$ in the group $(\mathbb{Z}/n\mathbb{Z})^{*}$ (Lesson 1's group again): it says $e^{k+1}\equiv e\pmod{\mathrm{ord}_n(m)}$. Recall $e$ was chosen coprime to $\varphi(n)$ during key generation, and $\mathrm{ord}_n(m)$ always divides $\varphi(n)$ (an element's order divides the group's size — this is the same Lagrange-style fact behind Euler's theorem), so $e$ is automatically coprime to $\mathrm{ord}_n(m)$ too, and therefore invertible modulo $\mathrm{ord}_n(m)$. Multiplying both sides of $e^{k+1}\equiv e$ by $e^{-1}\bmod\mathrm{ord}_n(m)$ cancels one factor of $e$ cleanly:
\[
e^{k} \equiv 1 \pmod{\mathrm{ord}_n(m)}.
\]

\textbf{Step 4: read off the plaintext.} Substituting this back, the term one step before the cycle closes is
\[
c_{k-1} \equiv m^{e^{k}} \equiv m^{1} \equiv m \pmod n,
\]
using Step 3 exactly as Lesson 1 used Euler's theorem: raising to an exponent that is $\equiv1$ modulo the element's order leaves the element unchanged.
\end{proof}

\medskip\noindent\textbf{The one-sentence summary.} Repeated encryption is iteration inside a finite cyclic structure — it \emph{must} eventually loop back to the starting ciphertext, and by the same ``one full lap plus one step'' logic from Lesson 1's Euler's-theorem walkthrough, the plaintext is sitting exactly one step before the loop closes.

\begin{example}[Cycling attack, fully worked]
Let $p=47, q=59$, so $n=2773$, $\varphi(n)=2668$, and let $e=17$. A message $m=123$ is encrypted: $c = 123^{17} \bmod 2773 = 1924$.

Repeatedly raising to the $e$-th power: $c_1 = c^{17} \bmod n, c_2 = c_1^{17} \bmod n, \dots$ — after $\mathbf{44}$ iterations, the sequence returns to $c_{44} = 1924 = c$. The value one step before, $c_{43} = 123 = m$, is the recovered plaintext.
\end{example}

\medskip\noindent\textbf{Why this stays purely conceptual.} The cycle length here is the multiplicative order of $e$ in $(\mathbb{Z}/\lambda(n)\mathbb{Z})^{*}$, which for well-chosen keys (large random primes, standard $e$) is astronomically large — comparable to $\varphi(n)$ itself in the worst case for the attacker. At real key sizes this attack is utterly infeasible; it survives in courses because it cleanly demonstrates that RSA encryption is literally iteration inside a finite cyclic structure, and it foreshadows why key-generation standards historically recommended checking that $\lambda(n)$ has large prime factors (to keep every element's order large) — the same underlying concern as smooth $p-1$ in Pollard's method (Lesson 1).

\newpage
\section{The Franklin--Reiter Related-Message Attack}

\subsection{Setup}
Suppose the \emph{same recipient} (same $n, e$) receives two ciphertexts of two messages known to satisfy a \textbf{publicly known linear relationship} $m_2 = m_1 + t$ for some known constant $t$ (e.g., two near-identical form-letter messages differing only in a numeric field), both under a small exponent, e.g. $e=3$:
\[
c_1 = m_1^{3} \bmod n, \qquad c_2 = (m_1+t)^{3} \bmod n.
\]

\medskip\noindent\textbf{Why this should feel familiar, before the proof.} Every earlier attack in this lecture combined two \emph{numbers} via the Euclidean algorithm (Wiener: numerator/denominator pairs; common modulus: two exponents) to isolate an unknown. Franklin--Reiter does the same trick one level of abstraction up: instead of combining two \emph{integers}, it combines two \emph{polynomials} — and the ``Euclidean algorithm'' for polynomials (repeated division with remainder, exactly as with integers, just tracking coefficients instead of digits) plays exactly the same role. The key structural fact is that both polynomials constructed below happen to share $m_1$ as a root purely by how they were built from public information — so their greatest common divisor must contain the factor $(x-m_1)$, and generically nothing else survives.

\begin{theorem}[Franklin--Reiter, 1996]
Given $(n, e{=}3, t, c_1, c_2)$, $m_1$ can be recovered in polynomial time, without factoring $n$.
\end{theorem}

\begin{proof}[Proof, step by step]
\textbf{Step 1: build two polynomials that are both public knowledge to construct, and both have $m_1$ as a root.} Define, over $\mathbb{Z}_n[x]$ (polynomials with coefficients reduced mod $n$ — everything needed to write these down is public):
\[
f_1(x) = x^3 - c_1, \qquad f_2(x) = (x+t)^3 - c_2.
\]
Substituting $x=m_1$ into $f_1$ gives $m_1^3-c_1=0$ by the definition of $c_1$; substituting $x=m_1$ into $f_2$ gives $(m_1+t)^3-c_2=0$ by the definition of $c_2$. So $x=m_1$ is a genuine root of \emph{both} polynomials, even though the attacker does not know $m_1$ — only that it exists as a common root.

\textbf{Step 2: a shared root means $(x-m_1)$ divides the gcd.} Basic polynomial algebra: if $x=m_1$ is a root of a polynomial $f$, then $(x-m_1)$ divides $f(x)$ exactly (polynomial division by a linear factor with zero remainder — the polynomial analogue of ``if $a$ divides $b$ with remainder $0$''). Since $(x-m_1)$ divides both $f_1$ and $f_2$, it divides any combination the Euclidean algorithm produces from them, and in particular it divides $\gcd(f_1,f_2)$.

\textbf{Step 3: argue the gcd is \emph{exactly} $(x-m_1)$, generically.} $f_1$ has degree $3$ and $f_2$ has degree $3$; their gcd could, in principle, be larger than the single factor $(x-m_1)$ if the two cubics happened to share some \emph{other} root too. But for two cubics built from unrelated random-looking ciphertexts, sharing a second common root ``by accident'' happens with negligible probability — so in essentially every real instance, $(x-m_1)$ \emph{is} the entire gcd, up to an irrelevant constant multiple. A linear polynomial's root is read off immediately.

\textbf{Step 4: compute the gcd mechanically, via the polynomial Euclidean algorithm.} Exactly as with integers — repeatedly divide the higher-degree polynomial by the lower-degree one and keep the remainder, until the remainder's degree can drop no further — except now every arithmetic operation on coefficients (add, subtract, multiply, and crucially \emph{divide}, which needs a modular inverse) is carried out modulo $n$. The process terminates at a linear polynomial $\alpha x+\beta\equiv0\pmod n$, whose unique root $x\equiv-\beta\alpha^{-1}\pmod n$ is $m_1$.
\end{proof}

\medskip\noindent\textbf{The one-sentence summary.} Two polynomials, built entirely from public data, are engineered by construction to share exactly one common root — the secret message — and the same Euclidean-algorithm machinery that finds integer gcds finds polynomial gcds just as mechanically, extracting that shared root directly.

\begin{example}[Franklin--Reiter, fully worked]
Let $n=3233$, $e=3$, and suppose $m_1 = 20$, with a second message $m_2 = m_1 + t = 25$ for known offset $t=5$. Ciphertexts: $c_1 = 20^3 \bmod 3233 = 1534$, $c_2 = 25^3 \bmod 3233 = 2693$.

Build $f_1(x) = x^3 - 1534$ and $f_2(x) = (x+5)^3 - 2693 = x^3 + 15x^2 + 75x + (125-2693)$ over $\mathbb{Z}_{3233}[x]$. Running the polynomial Euclidean algorithm (dividing $f_1$ by $f_2$, then continuing with remainders, all coefficient arithmetic done mod $3233$) reduces the pair to a single \textbf{linear} polynomial: $525x + 2432 \equiv 0 \pmod{3233}$. Solving: $x \equiv -2432 \cdot 525^{-1} \equiv 20 \pmod{3233}$ — exactly $m_1$. \textbf{Neither $n$ was factored nor any private key touched} — only public ciphertexts, the known offset $t$, and polynomial arithmetic.
\end{example}

\medskip\noindent\textbf{Practical consequence.} \textbf{Never} encrypt structurally related plaintexts (form letters, sequential IDs, incrementing counters) under a small exponent without randomized padding. This is one of the strongest practical arguments for OAEP: it destroys any algebraic relationship between the padded plaintexts, even when the underlying messages are related, which is exactly what this attack (and its harder unknown-relation cousin below) depends on.

\newpage
\section{Coppersmith's Method: An Overview}\label{sec:coppersmith}

\subsection{The General Problem}

\medskip\noindent\textbf{What problem is being solved, and why it is hard in general.} Every attack so far in this lecture found an unknown by exploiting some special structure of the specific scenario (small $d$, shared modulus, related messages). Coppersmith's method attacks a much more general question directly: \emph{given a polynomial equation that holds modulo $n$, when can we find a root, even without knowing how to factor $n$?} Solving polynomial equations modulo a composite number of unknown factorization is, in general, believed to be hard — it is essentially as hard as factoring $n$ itself for a typical polynomial. Coppersmith's remarkable result is that this hardness has an exception precisely when the root you are looking for is \emph{unusually small} relative to $n$ and the polynomial's degree.

\medskip\noindent\textbf{Coppersmith's small-root theorem (informal statement).} Given a polynomial $f(x)$ of degree $\delta$ modulo $n$ (with $n$ of unknown or partially unknown factorization), if $f$ has a root $x_0$ with $|x_0| < n^{1/\delta}$, then $x_0$ can be found in time polynomial in $\log n$ and $\delta$.

\medskip\noindent\textbf{Why smallness is exactly the right lever.} Think of it this way: an arbitrary root modulo $n$ carries roughly $\log_2 n$ bits of information — as much as $n$ itself — so pinning it down should be exactly as hard as factoring. But a root promised to satisfy $|x_0|<n^{1/\delta}$ carries only about $\frac{1}{\delta}\log_2 n$ bits — a small fraction of a full-size unknown. Coppersmith's theorem says that whenever the \emph{unknown part} of the problem is this small, there is enough redundancy in the surrounding structure (made precise by the lattice construction below) to pin it down efficiently, without ever needing to factor $n$ along the way.

This single theorem — due to Coppersmith (1996), with a simplified proof by Howgrave-Graham — underlies an entire family of RSA attacks that all reduce to ``find a small root of a polynomial modulo $n$'':

\begin{itemize}
\item \textbf{Håstad's short-pad attack} (the unknown-relation generalization of Håstad's broadcast attack above): recovers $m$ when two related messages are sent under low $e$ even when the \emph{relationship itself} is not fully known, only that it is small.
\item \textbf{Partial key exposure}: given enough of the high-order (or low-order) bits of $p$, recover the rest of $p$ and factor $n$ — directly relevant to side-channel and cold-boot attacks where an attacker recovers a fraction of the key bits.
\item \textbf{Boneh--Durfee}: an improvement on Wiener's bound, extending the recoverable range of small $d$ from $n^{0.25}$ up to $n^{0.292}$, using a two-variable version of Coppersmith's technique.
\item \textbf{Stereotyped-message attacks}: recovering $m$ when only a small unknown portion of a structured plaintext (e.g., a fixed-format message with one unknown short field) is unknown.
\end{itemize}

\medskip\noindent\textbf{The common shape of all four.} Look closely and every bullet above is the same sentence with different nouns substituted in: ``most of some quantity is known; a small remaining unknown piece satisfies a polynomial relationship mod $n$; recover the small piece.'' Recognizing a new attack scenario as an instance of this one shape is the single most useful skill Coppersmith's method teaches — the hard mathematical work (below) only has to be understood once.

\subsection{The Reduction to Lattices}

\medskip\noindent\textbf{The high-level plan, before the machinery.} A direct search for $x_0$ modulo $n$ is hard because $n$'s factorization is unknown — ordinary polynomial root-finding techniques need to work mod a prime or over the reals, not mod a mystery composite. Coppersmith's trick sidesteps this entirely: instead of solving the polynomial equation modulo $n$ directly, build \emph{several} polynomials that are all guaranteed to share the root $x_0$ modulo $n$, and search for a clever \emph{integer} combination of them that is small enough to force the same equation to hold as an honest equation over the integers $\mathbb{Z}$ — a domain where root-finding is easy and needs no knowledge of $n$'s factors at all. ``Searching for a clever small combination'' is precisely the Shortest Vector Problem in disguise, which is why lattice reduction enters the picture.

\medskip\noindent\textbf{From polynomial roots to short vectors.} The key trick: build several polynomials $g_1(x), g_2(x), \dots, g_k(x)$, all sharing the same root $x_0 \bmod n$ (typically shifts and powers of the original $f$, multiplied by powers of $n$ to control the modulus), and treat their \emph{coefficient vectors} as a basis for a lattice $L \subset \mathbb{Z}^k$. A theorem of Howgrave-Graham says: if $h(x) = \sum h_i x^i$ is an integer linear combination of the $g_i$ (hence still has $x_0$ as a root mod $n$) with \textbf{small enough Euclidean norm} $\|h\| < n/\sqrt{k}$, then $h(x_0) = 0$ holds as an equation over $\mathbb{Z}$, \emph{not just modulo $n$}. Since $h$ is now known to vanish at $x_0$ as an ordinary integer polynomial, standard root-finding recovers $x_0$ directly.

\medskip\noindent\textbf{Why the norm bound is exactly the right threshold — the intuition behind Howgrave-Graham's lemma.} If $h(x_0)\equiv0\pmod n$, then $h(x_0)$ is \emph{some} multiple of $n$ — possibly $0$, possibly $n$, possibly $-2n$, and so on. If we can additionally guarantee $h(x_0)$ is smaller in absolute value than $n$ itself, the only multiple of $n$ small enough to fit is $0$ — forcing $h(x_0)=0$ exactly. Bounding $|h(x_0)|$ in terms of $h$'s coefficient sizes (via the Euclidean norm $\|h\|$) and the known bound on $x_0$ is exactly what turns ``small coefficients'' into ``small value at $x_0$'', which is why finding a short vector in the lattice $L$ is the right computational target.

So the whole problem reduces to: \textbf{find a short vector in the lattice $L$} spanned by the $g_i$'s coefficient vectors. This is precisely Shortest Vector Problem territory — and in practice, the LLL algorithm (an approximate-SVP solver) is exactly the tool used, since Coppersmith's bound is calibrated to what LLL can reliably deliver in polynomial time. If you have covered lattice-based cryptography, this is the same LLL you met there, now redeployed as an attack tool rather than a hardness-assumption underpinning.

\begin{example}[Toy illustration of the core mechanism — small-scale, not full LLL]
Suppose we know $p$ shares its top bits with a known guess $\tilde{p}$: specifically $p = \tilde p + x_0$ for some small unknown $x_0$, and we are told $|x_0| < n^{1/2}$ (deliberately generous, purely for illustration). Consider $f(x) = \tilde p + x$, whose root mod $p$ is $x_0$ by construction. In the real attack, one builds a small lattice from shifted copies of $f$, powers of $n$ or $p$, and known bounds on $x_0$, then runs LLL to extract a short combination $h(x)$ satisfying Howgrave-Graham's norm bound; solving $h(x_0)=0$ over $\mathbb{Z}$ (an ordinary, easy root-finding step once $h$ is small) recovers $x_0$, hence $p$, hence the full factorization of $n$.

At the tiny numeric scale used throughout this lecture (3–4 digit moduli), the lattices involved collapse to dimension 1 or 2 and the bound on $x_0$ needed for a compelling illustration is not usefully separable from simple brute-force search — so a genuinely small worked numeric example does not illustrate the mechanism honestly. This is a case where \textbf{the technique's value only appears at realistic scale} (hundreds of bits), which is exactly why Coppersmith's method is taught as a reduction to a problem you already have a tool for (LLL), rather than by hand-computation as in the earlier attacks in this lecture.
\end{example}

\begin{remark}[Why this belongs at the end of this lecture, not the start]
Every attack earlier in this lecture (Wiener, Håstad, common modulus, Franklin--Reiter) can be fully hand-verified with small numbers, which is why they came first. Coppersmith's method is the most powerful of the family but the least hand-computable — its correctness rests on LLL actually finding a short-enough vector, which is a statement about lattice reduction quality, not elementary number theory. This is intentionally the conceptual capstone of RSA's mathematical attacks, and it is the point where number theory hands off to lattice reduction (the LLL algorithm) as the working tool.
\end{remark}

\newpage
\section{Summary and Bridge to Lesson 3}

\subsection*{Summary}
\begin{itemize}
\item \textbf{Wiener's attack}: a small $d$ is recoverable directly from $(n,e)$ via continued fractions — never choose $d$ for efficiency.
\item \textbf{Håstad's broadcast attack}: low exponent + repeated/related messages to multiple recipients + no padding = full break via CRT.
\item \textbf{Common modulus attack}: sharing $n$ across users with coprime exponents is a full break via the Extended Euclidean Algorithm.
\item \textbf{Cycling attack}: a conceptual reminder that RSA encryption is iteration in a finite group — impractical, but foundational.
\item \textbf{Franklin--Reiter}: known-relation messages under low exponent break via polynomial GCD.
\item \textbf{Coppersmith's method}: the general small-root machinery underlying stronger versions of several of the above, reducing to LLL lattice reduction.
\end{itemize}
\textbf{The unifying lesson}: every attack here targets a specific \emph{misuse} of RSA — small exponents without padding, small private exponents, repeated moduli, related plaintexts. None of them factor $n$, and all of them are defeated by proper randomized padding (OAEP) and standard key-generation discipline (large random $d$, non-shared $n$, large $e$ such as 65537).

\subsection*{Looking Ahead --- Lesson 3 — Implementation and Protocol Attacks on RSA}
Next lecture assumes RSA is used \emph{correctly} at the mathematical level (proper key generation, no shared moduli, no related-message broadcasting) and instead breaks it through its physical or protocol-level implementation:
\begin{itemize}
\item \textbf{Bleichenbacher's padding oracle attack} on PKCS\#1 v1.5 — recovering plaintext by observing only whether a server accepts or rejects a padding format.
\item \textbf{Timing attacks} exploiting the data-dependent branching in square-and-multiply (flagged back in Lesson 1).
\item \textbf{Fault attacks} on CRT-RSA (the Bellcore attack) — a single computational glitch during the $m_p$ or $m_q$ branch (Lesson 1, Section~\ref{sec:crt-speedup}) that factors $n$ from one faulty signature.
\item \textbf{ROCA} and other real-world key-generation vulnerabilities.
\item OAEP's security argument, and why it defeats every low-exponent attack in today's lecture.
\end{itemize}

\section*{Tutorial Exercises}
\begin{enumerate}
\item Run Wiener's attack by hand on $n = 100822548703$ with $e = 74531608157$ — hint: the continued fraction expansion has fewer than 10 terms. (Verify your recovered $d$ by checking $ed \equiv 1$ modulo your recovered $\varphi(n)$.)
\item In the Håstad broadcast example, what is the minimum number of recipients needed if $e = 5$ instead of $3$? Explain why in terms of the CRT bound $m^e < N$.
\item Extend the common modulus attack: what happens if $\gcd(e_1,e_2) = g > 1$? Is the attack still possible, and if so what can be recovered?
\item In the cycling attack, express the cycle length $k$ precisely in terms of the multiplicative order of $e$ modulo $\lambda(n)$ (the Carmichael function), and explain why choosing $p,q$ such that $\lambda(n)$ has a large prime factor defends against short cycles.
\item Redo the Franklin--Reiter worked example with $e=3$, $t=7$, $m_1 = 30$ on the same $n=3233$. Verify your recovered $m_1$.
\item Research question: read the original Boneh--Durfee paper's abstract and result statement. What is the improved bound on $d$ compared to Wiener's $n^{0.25}$, and what extra mathematical tool (beyond one-dimensional continued fractions) does it require?
\end{enumerate}


\chapter{Lesson 3: The RSA Cryptosystem III --- Implementation and Protocol Attacks}
\begin{center}
\large\itshape Padding Oracles, Timing and Fault Attacks, ROCA, OAEP, and the Guillou--Quisquater ZKP Bridge
\end{center}
\vspace{0.4cm}
\vspace{0.5cm}
\subsection*{Roadmap}
\begin{enumerate}[label=(\arabic*)]
\item Padding oracles: PKCS\#1 v1.5 and Bleichenbacher's ``million-message'' attack
\item Timing attacks: leaking secret exponent bits through execution time, and blinding as the fix
\item Fault attacks: the Bellcore attack on CRT-RSA signatures
\item ROCA: a real-world key-generation flaw that turned Coppersmith's method (Lesson 2) into a practical factoring attack
\item Side-channel countermeasures: a consolidated checklist
\item OAEP: a security sketch for the padding scheme that defeats almost everything in Lessons 2--3
\item Guillou--Quisquater: an RSA-based zero-knowledge proof of identity, and why it needs only one exponentiation-sized challenge space instead of many rounds
\end{enumerate}
\textbf{What every attack in this lecture has in common:} $n$ is properly generated, $d$ is large, no modulus is shared, no message is broadcast — every mathematical precondition from Lessons 1--2 is satisfied. These attacks instead exploit \emph{how RSA is computed} (timing, faults) or \emph{how a protocol reacts to malformed input} (padding oracles), or a subtle flaw in \emph{how the keys themselves were generated} (ROCA). This is the layer of RSA security that pure number theory cannot see.

\newpage
\section{Padding Oracles: Bleichenbacher's Attack on PKCS\#1 v1.5}

\subsection{PKCS\#1 v1.5 Encryption Padding}

Recall from Lesson 1 that textbook RSA is deterministic and malleable, and that real systems pad the message before encrypting. The older padding standard, \textbf{PKCS\#1 v1.5}, formats a $k$-byte block (where $k$ is the byte-length of $n$) as
\[
\texttt{EB} = \texttt{0x00} \,\|\, \texttt{0x02} \,\|\, \text{PS} \,\|\, \texttt{0x00} \,\|\, D,
\]
where $D$ is the message, $\text{PS}$ is a string of \emph{non-zero} random padding bytes at least 8 bytes long (padding out the block to exactly $k$ bytes), and the leading \texttt{0x00 0x02} is a fixed, publicly-known format marker. The recipient decrypts $c$, obtains $\text{EB} = c^d \bmod n$ (represented as a $k$-byte string), and checks the format: does it start with \texttt{0x00 0x02}, is there a \texttt{0x00} byte terminating the padding within a valid range, etc. If the format is invalid, many real implementations (TLS servers being the classic case) return a distinguishable error — \emph{this is the oracle}.

\begin{definition}[Padding oracle]
A \textbf{padding oracle} for PKCS\#1 v1.5 is any mechanism (an error message, a timing difference, a connection reset — the specific channel does not matter) that lets an attacker submit a chosen ciphertext $c$ and learn a single bit: \emph{does $c^d \bmod n$ decrypt to a validly-PKCS\#1-padded block, i.e.\ does it start with \texttt{0x00 0x02}?}
\end{definition}

\begin{example}[A tiny PKCS\#1 v1.5 block, made concrete]
Take a toy $k=8$-byte block (a real RSA-1024 block would use $k=128$; 8 bytes is enough to see the structure). Encoding the 2-byte message $D = \texttt{0x1234}$ with 4 bytes of random non-zero padding $\text{PS} = \texttt{0xA1 0x7F 0x3C 0x99}$ gives
\[
\texttt{EB} = \underbrace{\texttt{0x00}}_{\text{1 byte}} \,\underbrace{\texttt{0x02}}_{\text{1 byte}} \,\underbrace{\texttt{0xA1 0x7F 0x3C 0x99}}_{\text{PS, 4 bytes}} \,\underbrace{\texttt{0x00}}_{\text{1 byte}} \,\underbrace{\texttt{0x1234}}_{\text{message, 2 bytes}},
\]
exactly 8 bytes. A recipient decrypting some ciphertext $c'$ and obtaining, say, $\texttt{0x00 0x02 0x00 \dots}$ (a zero byte immediately after the \texttt{0x02} marker, where PS must be non-zero) rejects it as invalid — this is precisely the format check the padding oracle exposes as a single accept/reject bit, with nothing about the message $D$ itself revealed directly by the check.
\end{example}

\medskip\noindent\textbf{Why one bit is catastrophic.} A one-bit oracle sounds harmless. Bleichenbacher (CRYPTO 1998) showed it is not: an attacker who can query the oracle on ciphertexts of the form $c' = c \cdot s^e \bmod n$ for chosen \emph{blinding} values $s$ learns, for each query, whether $m \cdot s \bmod n$ lands in the narrow range $[2B, 3B)$ where $B = 2^{8(k-2)}$ (i.e., whether the padded value ``looks like'' \texttt{0x00 0x02...}). Each YES answer is a strong constraint on the unknown $m$. Repeating this adaptively — using each new constraint to choose better values of $s$ — narrows the set of possible $m$ until only one remains. This is a classic \textbf{adaptive chosen-ciphertext attack}: it uses the decryption oracle itself, not any weakness in the RSA problem.

\subsection{The Attack, at a High Level}

\begin{algorithm}[h]
\caption{Bleichenbacher's Attack (structure, not full case analysis)}
\begin{algorithmic}[1]
\State \textbf{Input:} target ciphertext $c = m^e \bmod n$; oracle $\mathrm{Valid}(c') \in \{0,1\}$ testing PKCS\#1 conformance
\State \textbf{Step 1 (blinding search):} find the smallest $s_1$ such that $c \cdot s_1^{e} \bmod n$ passes the oracle
\State \textbf{Step 2 (interval narrowing):} maintain a shrinking set of intervals that must contain $m$; each oracle query on a candidate $s_i$ either confirms an interval or splits it further, using the known structure $[2B,3B)$
\State \textbf{Step 3 (interval refinement):} once a single interval remains, search increasingly precise values of $s$ to shrink the interval toward a single point
\State \textbf{Return} the unique $m$ remaining once the interval has width 1
\end{algorithmic}
\end{algorithm}

\begin{remark}[Why ``million-message attack'']
Bleichenbacher's original paper estimated roughly $2^{20}$ (about a million) oracle queries to fully decrypt one ciphertext against 1024-bit RSA-PKCS\#1v1.5 — expensive but entirely practical for an attacker with network access to a server that leaks the oracle bit, e.g.\ via a distinguishable TLS alert message (``bad padding'' vs.\ other decryption failures). Later refinements (Bardou et al., 2012; ROBOT, 2017) reduced the required query count and revealed the attack was still exploitable in production TLS stacks nearly 20 years after publication — the flaw is in the \emph{error-handling contract} of the protocol, which is easy to get subtly wrong even when the padding format itself is implemented correctly.
\end{remark}

\medskip\noindent\textbf{Practical consequence.} \textbf{Never} let decryption failure reasons be distinguishable to the attacker — not via error message content, not via response timing, not via connection behavior. The historically recommended mitigation for legacy PKCS\#1 v1.5 is to proceed \emph{as if} decryption succeeded on failure (substituting random plaintext) so that no observable signal exists, rather than trying to enumerate and suppress every possible error channel. The complete fix, however, is to stop using PKCS\#1 v1.5 for new designs and use OAEP (Section~\ref{sec:oaep}), whose security proof directly rules out this class of attack.

\newpage
\section{Timing Attacks}

\subsection{The Leakage Source}

Recall the square-and-multiply algorithm from Lesson 1: the \texttt{if} branch executes an extra multiplication exactly when the current secret exponent bit is $1$. Even without a visible branch, the multiplication and squaring operations themselves can take data-dependent time (due to conditional reduction steps, cache access patterns, or hardware multiplier behavior on operands of different sizes/structure).

\medskip\noindent\textbf{Kocher's Timing Attack (1996), core idea.} An attacker who can submit many ciphertexts $c_1, c_2, \dots$ to a server and measure the wall-clock time $T_i$ each decryption/signature takes can recover the secret exponent $d$ \textbf{one bit at a time}, from the most significant bit downward. Having correctly guessed the top $j$ bits of $d$, the attacker can \emph{simulate} the modular exponentiation up to that point for each $c_i$ and correlate the simulated intermediate computation's cost with the real measured $T_i$. A correct bit guess produces a statistically detectable correlation (the simulated step matches what the real hardware actually did); a wrong guess does not. Repeating bit by bit, with enough samples per bit to beat measurement noise, recovers all of $d$.

\begin{remark}[Why this generalizes well beyond the toy square-and-multiply]
The same statistical idea applies to \emph{any} implementation where some intermediate computational step's cost depends on secret bits — including CRT-RSA's $d_p, d_q$ exponentiations (arguably a richer target, since each exponent is only half the length, so fewer bits need to be recovered per correlated measurement), and even to cache-timing variants that measure which memory lines were touched rather than wall-clock time directly (e.g., a colocated attacker on the same physical machine observing cache eviction patterns).
\end{remark}

\begin{example}[Toy illustration of the bit-by-bit correlation]
Let the secret exponent be the 6-bit toy value $d = 101101_2 = 45$, used with right-to-left square-and-multiply (Lesson 1, Algorithm~\ref{alg:modexp}). Processing $d$'s bits from the \emph{least}-significant end, the algorithm performs a squaring every step, plus an \emph{extra} multiplication exactly on the steps where the bit is $1$: reading $d=101101_2$ least-significant-bit first gives the bit sequence $1,0,1,1,0,1$, so extra multiplications occur on steps $1,3,4,6$ and are skipped on steps $2,5$. An attacker who has correctly guessed the first few (low-order) bits can simulate exactly which steps \emph{should} include the extra multiplication; timing many real executions and checking whether the measured cost is high (extra multiply present) or low (skipped) on the predicted step is precisely the one-bit correlation test Kocher's attack performs, repeated bit by bit until all 6 bits of $d$ are recovered. At real key sizes ($d$ having thousands of bits, not 6), each bit requires many timing samples to beat measurement noise, but the mechanism is identical to this toy trace.
\end{example}

\subsection{Defense: Blinding}

\begin{definition}[RSA blinding]
Before computing $m = c^d \bmod n$, the signer/decryptor picks a fresh random $r$ (coprime to $n$) for every operation, computes
\[
c' = c \cdot r^{e} \bmod n, \qquad m' = (c')^{d} \bmod n, \qquad m = m' \cdot r^{-1} \bmod n.
\]
\end{definition}

\begin{proposition}[Correctness of blinding]
$m' = (c \cdot r^e)^d = c^d \cdot r^{ed} \equiv c^d \cdot r \pmod n$ (using $r^{ed} \equiv r \pmod n$ by the same Euler's-theorem identity from Lesson 1 that makes RSA correct in the first place), so $m' \cdot r^{-1} \equiv c^d \equiv m \pmod n$.
\end{proposition}

\medskip\noindent\textbf{Why blinding defeats the timing attack.} The exponentiation the attacker times is now $(c')^d \bmod n$ for a value $c'$ that the attacker cannot predict or control (it is masked by a fresh random $r$ chosen internally, unknown to the attacker, on every single call). Since the attacker cannot correlate the timing of a specific known input against a specific bit-guess simulation — the actual base being exponentiated is different and unknown every time — the statistical correlation Kocher's attack relies on collapses to noise. \textbf{Blinding does not make the exponentiation itself constant-time}; it removes the attacker's ability to choose or know the input being timed, which is what the attack actually needs.

\begin{example}[Blinding, fully worked]
Reuse the familiar toy key $p=61,q=53,n=3233,e=17,d=2753$. An attacker sends ciphertext $c = 65^{17} \bmod 3233 = 2790$ (encryption of $m=65$) and wants to time the server's decryption of exactly this $c$.

The server instead picks a fresh random blinding factor $r=41$ (coprime to $n$) and computes
\[
c' = c \cdot r^{17} \bmod 3233 = 2130.
\]
It is $c'=2130$ — not the attacker's original $c=2790$ — whose decryption the server actually times: $m' = (c')^{d} \bmod n = 2130^{2753} \bmod 3233 = 2665$. Unblinding recovers the true answer:
\[
m = m' \cdot r^{-1} \bmod n = 2665 \cdot 41^{-1} \bmod 3233 = 65,
\]
correctly matching the original message. \textbf{The attacker's timing measurement is on $c'=2130$, a value they never chose and cannot predict} — a fresh random $r$ is drawn on every call, so the same $c$ is never timed twice under the same disguise, and no bit-guess simulation the attacker builds around their chosen $c$ can be correlated against what the server actually computed.
\end{example}

\medskip\noindent\textbf{Practical consequence.} Production RSA implementations combine \emph{two} defenses, not one: (1) blinding, to remove attacker control over the timed input, and (2) constant-time modular exponentiation (e.g., Montgomery ladder instead of naive square-and-multiply, avoiding secret-dependent branches or table lookups entirely) as defense in depth, since blinding alone does not stop attackers who can observe \emph{which physical operations occurred} rather than just aggregate time (e.g., some cache/microarchitectural attacks). Real-world CVEs against RSA timing side channels (e.g., against OpenSSL's Montgomery multiplication in various years) have repeatedly targeted implementations that had blinding but incomplete constant-time arithmetic underneath it.

\newpage
\section{Fault Attacks: The Bellcore Attack on CRT-RSA}

\subsection{Setup}

Recall CRT-RSA decryption/signing from Lesson 1 (Section~\ref{sec:crt-speedup}): the signer computes
\[
s_p = c^{d_p} \bmod p, \qquad s_q = c^{d_q} \bmod q,
\]
and recombines via CRT to get the full signature $s$. This is roughly $4\times$ faster than computing $c^d \bmod n$ directly — which is precisely why real implementations use it, and precisely what makes the following attack possible.

\begin{theorem}[Bellcore attack, Boneh--DeMillo--Lipton, 1997]
Suppose an attacker can induce a single transient computational fault (via voltage glitching, clock glitching, laser fault injection, or simply naturally-occurring hardware error) during \emph{exactly one} of the two CRT branches — say the $s_p$ computation — while the $s_q$ branch computes correctly, and can obtain the resulting faulty signature $s'$ on a message $m$ whose value the attacker knows. Then the attacker recovers the prime factor $q$ (and hence the full private key) in polynomial time, using only $s'$, $m$, and the public $(n,e)$ — \textbf{no correct signature is needed at all.}
\end{theorem}

\begin{proof}[Proof idea]
Since only the $s_p$ branch was corrupted, the faulty $s'$ still satisfies $s' \equiv m^d \pmod q$ (the $q$-branch was untouched), but $s' \not\equiv m^d \pmod p$ (the $p$-branch was corrupted, and a random fault essentially never happens to still satisfy the same congruence). Hence
\[
q \mid (s'^{e} - m), \qquad p \nmid (s'^{e} - m)
\]
(the second holds because $s'^e \equiv m \pmod q$ requires $s' \equiv m^d \pmod q$ by RSA correctness restricted mod $q$, which holds, but the same argument mod $p$ fails since $s'$ is wrong there). Therefore $\gcd(s'^e - m \bmod n,\ n) = q$ exactly — the Euclidean algorithm reveals the fault-affected prime directly, and from $q$ the attacker computes $p = n/q$, then the entire private key.
\end{proof}

\begin{example}[Bellcore attack, fully worked]
Using the familiar toy key $p=61, q=53, n=3233, e=17, d=2753$, with $d_p = 53, d_q = 49, q^{-1} \bmod p = 38$. The signer correctly computes $s_p = 65^{53} \bmod 61 = 4$ and $s_q = 65^{49} \bmod 53 = 12$, recombining to the correct signature $s = 588$ (indeed $588^{17} \bmod 3233 = 65 = m$, as required).

Now suppose a fault corrupts the $p$-branch only, flipping $s_p$ to $s_p' = 5$ (one off from the correct value) while $s_q = 12$ is untouched. Recombining via the same CRT formula gives a faulty signature $s' = 2602$. Checking: $(s')^{17} \bmod 3233 = 1496$, and indeed $1496 \equiv 12 \equiv m \pmod{53}$ (the $q$-branch congruence survives) while $1496 \equiv 32 \not\equiv 4 \equiv m \pmod{61}$ (the $p$-branch congruence is broken).

The attacker, knowing only $n=3233$, $e=17$, the message $m=65$, and the single faulty signature $s'=2602$, computes
\[
\gcd\big((s')^{e} \bmod n \;-\; m,\ n\big) = \gcd(1496 - 65,\ 3233) = \gcd(1431, 3233) = \mathbf{53} = q.
\]
\textbf{A single faulty signature factored the modulus completely} — recovering $q$, then $p = 3233/53 = 61$, then $\varphi(n)$ and $d$ itself. No correct signature was ever needed.
\end{example}

\medskip\noindent\textbf{Practical consequence: signature verify-before-release.} The standard defense is procedural, not algebraic: \textbf{always verify a freshly computed CRT-RSA signature against the public key before releasing it} — i.e.\ check $s^e \bmod n \stackrel{?}{=} m$ internally, and discard/retry silently if the check fails, releasing nothing. A genuine fault will (overwhelmingly likely) fail this check, so the faulty signature the attacker needs is simply never emitted. This single check closes the entire attack class at negligible extra computational cost (one public-exponent exponentiation, which is cheap since $e$ is small, e.g.\ 65537), and is mandated by essentially every modern smart-card and HSM signing standard.

\begin{remark}[Why this specifically targets CRT-RSA, and why that matters]
Non-CRT RSA (computing $c^d \bmod n$ directly) is \emph{also} vulnerable to fault attacks in principle, but recovering the key from a single faulty full-size exponentiation is far harder — there is no clean two-branch structure to exploit algebraically. This is precisely the trade-off flagged in Lesson 1: the $4\times$ CRT speedup is bought at the cost of a much sharper fault-attack surface, which is why smart cards and HSMs that use CRT-RSA \emph{must} implement the verify-before-release countermeasure (or an equivalent redundant computation), while implementations that skip CRT entirely face a smaller (but nonzero) fault-attack risk.
\end{remark}

\newpage
\section{ROCA: When Key Generation Itself Is the Vulnerability}

\subsection{Background}

\textbf{ROCA} (``Return of Coppersmith's Attack,'' Nemec et al., 2017) is not a flaw in RSA's mathematics or in a protocol's error handling — it is a flaw in a specific \emph{key-generation implementation} (the Infineon RSALib used in many smart cards, TPMs, and government-issued ID cards) that made factoring \emph{practically feasible} for affected keys, at sizes (512--2048 bits) that are otherwise completely safe against GNFS.

\medskip\noindent\textbf{The vulnerable construction.} The affected library generated primes of the special form
\[
p = k \cdot M + \big(65537^{a} \bmod M\big)
\]
for a secret integer $k$, a secret exponent $a$, and a fixed public constant $M$ (a product of the first several hundred small primes — a ``primorial''), apparently as a performance optimization for the prime-generation process on constrained smart-card hardware. This construction was undisclosed until reverse-engineered by researchers.

\begin{proposition}[Why this collapses the search space]
An attacker who knows the public modulus $n = pq$ and suspects it was generated this way does \emph{not} need to search all $\approx \sqrt{n}$ possible values of $p$. Instead, since $p \bmod M$ is one of only $\mathrm{ord}_M(65537)$ possible residues (the multiplicative order of $65537$ modulo $M$ — typically a number in the thousands to low millions for the real ROCA parameters, versus $M$ itself and $\sqrt n$ being astronomically larger), the attacker can:
\begin{enumerate}
\item Enumerate the (comparatively tiny) set of candidate residues $r_a = 65537^a \bmod M$.
\item For each candidate residue, use \textbf{Coppersmith's method} (Lesson 2, Section~\ref{sec:coppersmith}) to search for a small $k$ completing $p = kM + r_a$ into an actual factor of $n$ — Coppersmith's small-root machinery is exactly suited to this ``most of $p$ is unknown but constrained to a coset'' structure.
\end{enumerate}
This turns a brute-force search over all of $\mathbb{Z}$ into a brute-force search over a small candidate set, \emph{composed with} a polynomial-time root-finding step per candidate — dropping the effective attack cost for affected 1024-bit keys from ``infeasible'' to a few CPU-hours, and for affected 2048-bit keys to a few CPU-months on modest hardware (the original paper's reported figures), rather than the millennia GNFS would require.
\end{proposition}

\begin{example}[Toy illustration of the fingerprint structure — not real ROCA scale]
Let the (deliberately tiny, illustrative-only) constant be $M = 3 \cdot 5 \cdot 7 \cdot 11 = 1155$. Compute $65537 \bmod M = 857$; the multiplicative order of $857$ modulo $1155$ turns out to be $12$ — meaning there are only \textbf{12} distinct residues $65537^a \bmod M$ ever reachable, out of the $1155$ residues mod $M$ in general (and utterly negligible next to $\sqrt n$ for a cryptographic-size $n$). Taking $a=1$ (so the residue is $65537 \bmod 1155 = 857$) and $k=2$ gives the candidate
\[
p = 2 \cdot 1155 + 857 = 3167,
\]
which is in fact prime. An attacker facing a modulus $n = p \cdot q$ built this way would not need to guess $p$ among all integers near $\sqrt n$ — only among the $12$ residues mod $1155$ (here, $\{1, 331, 362, 529, 593, 692, 694, 857, 923, 991, 1024, 1088\}$), each followed by a Coppersmith-style search for $k$. At real ROCA parameters ($M$ a $\sim$$2^{200}$-ish primorial, order in the low millions), the identical structural idea — not this exact toy arithmetic — is what makes the real attack practical.
\end{example}

\medskip\noindent\textbf{Practical consequence.} This vulnerability class does not indicate any weakness in RSA itself, in Coppersmith's method, or in any of the earlier lectures' key-generation advice (large random independent primes, avoiding smooth $p-1$, avoiding shared/related moduli) — it is precisely what happens when an implementation \emph{silently deviates} from that advice for an unstated performance reason. The lesson for practitioners: \textbf{never trust that ``a widely deployed, certified library'' generates primes uniformly at random} without independent verification; ROCA affected certified smart cards and government ID infrastructure for years before detection, precisely because the deviation was invisible from the public key alone unless someone specifically tested for the fingerprint.

\newpage
\section{Side-Channel Countermeasures: A Consolidated Checklist}

The individual defenses mentioned throughout this lecture form a coherent set of standard practices. Collecting them here as a reference:

\begin{center}
\begin{tabular}{p{4cm}p{6.5cm}p{3.5cm}}
\toprule
\textbf{Threat} & \textbf{Countermeasure} & \textbf{Covered in} \\
\midrule
Padding-oracle CCA & OAEP instead of PKCS\#1 v1.5; uniform error handling & Sec.~1, Sec.~6 \\
Timing leakage & RSA blinding (randomize the input) & Sec.~2 \\
Timing/cache leakage & Constant-time modular exponentiation (Montgomery ladder) & Sec.~2 \\
CRT fault injection & Verify signature before release; redundant computation & Sec.~3 \\
Weak key generation & Independent verification of prime-generation randomness; avoid undisclosed proprietary shortcuts & Sec.~4 \\
General implementation drift & Use vetted, widely-audited cryptographic libraries; do not hand-roll modular exponentiation or padding logic & All sections \\
\bottomrule
\end{tabular}
\end{center}

\begin{remark}[The unifying theme, one more time]
Every defense above targets the \emph{gap between the mathematical specification of RSA and its physical realization}. Lessons 1--2 taught you when RSA is mathematically sound; this lecture teaches you that mathematical soundness is necessary but nowhere near sufficient — an implementation can be simultaneously ``textbook-perfect'' in its algebra and completely broken in its execution.
\end{remark}

\newpage
\section{OAEP: A Security Sketch}\label{sec:oaep}

\subsection{Construction}

\textbf{Optimal Asymmetric Encryption Padding} (Bellare--Rogaway, 1994) replaces the fixed, structured PKCS\#1 v1.5 padding with a randomized, ``all-or-nothing''-looking transform, before the result is encrypted with textbook RSA:

\medskip\noindent\textbf{OAEP encoding (simplified).} To encode a message $m$ (padded with a fixed-length string of zero bits $\text{PS}$ for domain separation) using two hash functions $G$ (a mask-generating function) and $H$, and a fresh random seed $r$:
\[
X = (m \,\|\, \text{PS}) \oplus G(r), \qquad Y = r \oplus H(X), \qquad \text{EM} = X \,\|\, Y,
\]
then $\text{EM}$ (interpreted as an integer $< n$) is encrypted as ordinary textbook RSA: $c = \text{EM}^e \bmod n$. Decryption inverts each XOR step in reverse (recovering $r$ from $Y$ and $H(X)$, then $m \,\|\, \text{PS}$ from $X$ and $G(r)$) and \textbf{rejects} if the recovered padding string does not exactly match the expected $\text{PS}$.

\begin{example}[Toy illustration of the OAEP masking mechanism — not a real hash]
The real $G,H$ are cryptographic hash functions; to see the seed/mask mechanics with actual numbers, use toy 4-bit stand-ins $G(r) = (7r+3) \bmod 16$ and $H(X) = (5X+1) \bmod 16$ (fixed, public functions — not secret), omit $\text{PS}$ for brevity, and encode the toy message $m=5$ with fresh random seed $r=9$:
\[
G(9) = 2, \qquad X = m \oplus G(r) = 5 \oplus 2 = 7, \qquad H(7) = 4, \qquad Y = r \oplus H(X) = 9 \oplus 4 = 13.
\]
So $\text{EM} = (X,Y) = (7,13)$. Decoding reverses the steps in the opposite order: $r' = Y \oplus H(X) = 13 \oplus 4 = 9$, then $m' = X \oplus G(r') = 7 \oplus 2 = 5$ — correctly recovering $m$.

Now encrypt the \emph{same} message $m=5$ again with a fresh seed $r=2$: $G(2)=1$, so $X = 5 \oplus 1 = 4$; $H(4)=5$, so $Y = 2 \oplus 5 = 7$; giving $\text{EM}=(4,7)$ — completely different from $(7,13)$ despite encoding the identical message. \textbf{This is exactly the property that defeats Håstad and Franklin--Reiter}: two encodings of the same (or related) plaintext no longer share any visible algebraic structure once OAEP has randomized them, before RSA is even applied.
\end{example}

\medskip\noindent\textbf{Security sketch, in the random oracle model.} Modeling $G$ and $H$ as independent random oracles, Bellare and Rogaway showed OAEP-RSA achieves \textbf{IND-CCA2 security} (indistinguishability under adaptive chosen-ciphertext attack) whenever the underlying trapdoor permutation (here, RSA) is one-way. The proof idea: any adversary that can distinguish encryptions, or that submits a chosen ciphertext to a decryption oracle and gets useful information back, can be shown (via a simulation/extraction argument) to have queried the random oracle $H$ or $G$ on the specific value tied to the challenge ciphertext's underlying RSA pre-image — which lets a simulator extract an RSA inverse from the adversary's queries, contradicting the one-wayness of RSA. Informally: \textbf{any successful attack on OAEP-RSA can be converted into an algorithm that inverts plain RSA}, which we assume is hard.

\begin{remark}[Why OAEP specifically defeats the earlier attacks]
Each attack in Lessons 2--3 relied on some algebraic or structural regularity surviving into the encrypted value:
\begin{itemize}
\item \textbf{Håstad / Franklin--Reiter} (Lesson 2) needed related plaintexts to remain related after encoding. OAEP's random seed $r$ makes two encryptions of related (or even identical) messages look like uniformly random, unrelated strings before RSA is even applied — there is no longer any algebraic relationship for CRT or polynomial-GCD tricks to exploit.
\item \textbf{Bleichenbacher's padding oracle} (Section~1) exploited the fact that a large fraction of PKCS\#1v1.5-decrypted garbage still ``looks like'' a differently-shaped valid message with nonzero probability, and that validity itself was cheaply, deterministically checkable and distinguishable. OAEP's redundancy check (the $\text{PS}$ match) has negligible probability of accidentally validating on a random/incorrect decryption, and — critically — OAEP's \emph{security proof} accounts for a decryption oracle directly (IND-CCA2 is defined with exactly such an oracle available to the adversary), which PKCS\#1v1.5's original design did not.
\item \textbf{Malleability and determinism} (flagged already in Lesson 1) are both eliminated: the random seed $r$ makes encryption probabilistic, and the tightly-coupled XOR/hash structure means a modified ciphertext decrypts (if it decrypts at all) to something the redundancy check almost always rejects, rather than to a predictably related plaintext.
\end{itemize}
\end{remark}

\medskip\noindent\textbf{What OAEP does \emph{not} fix.} OAEP is a defense at the \emph{mathematical/protocol} layer. It does nothing by itself against Section~2's timing attacks or Section~3's fault attacks against the underlying modular exponentiation — those require blinding, constant-time arithmetic, and verify-before-release, independently of which padding scheme sits on top. Real-world security requires \emph{all} of these layered together; OAEP alone is necessary but not sufficient.

\newpage
\section{Guillou--Quisquater: An RSA-Based Zero-Knowledge Proof}

\subsection{Motivation: Proving Identity Without Revealing the Secret}

Everything so far in this RSA module has been about encryption or signing. Guillou--Quisquater (1988) uses the same algebraic structure — an RSA-type modulus and modular exponentiation — for a different goal: letting a prover convince a verifier \emph{``I know a secret $s$ tied to my claimed identity''} without revealing $s$, and without letting the verifier (or an eavesdropper) later impersonate the prover.

\medskip\noindent\textbf{Setup.} A trusted authority publishes an RSA-like modulus $n = pq$ and a public verification exponent $v$ (coprime to $\varphi(n)$, exactly as with an RSA public exponent). Each prover's identity is mapped, via a public redundancy/hash function, to a value $J$; the authority (who alone knows $\varphi(n)$, i.e.\ $p,q$) issues the prover a secret credential $s$ satisfying
\[
s^{v} \cdot J \equiv 1 \pmod n, \qquad \text{i.e.}\qquad s \equiv J^{-1/v} \pmod n,
\]
computable by the authority (who can take $v$-th roots mod $n$ using $\varphi(n)$) but not by anyone else, exactly as RSA decryption requires the private exponent.

\begin{algorithm}[h]
\caption{Guillou--Quisquater Identification Protocol (one round)}\label{alg:gq}
\begin{algorithmic}[1]
\State \textbf{Prover} picks random $r$ coprime to $n$, sends commitment $x = r^{v} \bmod n$
\State \textbf{Verifier} picks random challenge $c \in \{0, 1, \dots, v-1\}$, sends $c$
\State \textbf{Prover} computes and sends $y = r \cdot s^{c} \bmod n$
\State \textbf{Verifier} accepts iff $x \equiv y^{v} \cdot J^{c} \pmod n$
\end{algorithmic}
\end{algorithm}

\begin{proposition}[Completeness]
If the prover genuinely knows $s$ and follows the protocol, $y^v \cdot J^c = (r s^c)^v \cdot J^c = r^v \cdot s^{vc} \cdot J^c = x \cdot (s^v J)^c \equiv x \cdot 1^c \equiv x \pmod n$, so an honest verifier always accepts.
\end{proposition}

\begin{proposition}[Soundness, sketch]\label{prop:gq-soundness}
A cheating prover who does not know $s$ but wants to pass for a specific challenge value $c$ can indeed fake a valid-looking transcript for \emph{that one} value of $c$ (by choosing $y$ first, then setting $x = y^v J^c \bmod n$ to match) — but must commit to $x$ \emph{before} learning the verifier's actual random challenge. A cheating prover can therefore satisfy at most one of the $v$ possible challenge values with a given commitment $x$, so guessing the verifier's challenge correctly happens with probability exactly $1/v$. Choosing $v$ large (e.g.\ a value comparable in size to a standard RSA public exponent) drives the cheating probability negligibly low in a \emph{single round}, unlike Fiat--Shamir-style protocols built on a $v=2$ structure, which need many sequential rounds to reach the same soundness level.
\end{proposition}

\begin{remark}[Zero-knowledge, sketch]
A simulator that does not know $s$ can still produce transcripts $(x,c,y)$ indistinguishable from real ones \emph{without} interacting with a real prover: pick $c$ and $y$ uniformly at random first, then set $x = y^v J^c \bmod n$ (exactly the ``rewind and fake the commitment'' trick used in Proposition~2, but now run by the simulator rather than an attacker). Since $x$ is computed from uniformly random $y,c$, the resulting triple has exactly the same distribution as a genuine protocol run, which is the essence of the honest-verifier zero-knowledge argument: nothing in the transcript could not have been produced without knowledge of $s$, so the transcript itself leaks no information about $s$.
\end{remark}

\begin{example}[Guillou--Quisquater, fully worked]
Reuse the familiar toy modulus $n = 3233$ ($p=61,q=53$, $\varphi(n)=3120$), and take the public exponent $v = 7$ (check: $\gcd(7,3120)=1$, exactly the RSA-key-generation condition from Lesson 1). Suppose the authority has issued the prover secret $s = 100$ (with $\gcd(100,3233)=1$), corresponding to public identity value
\[
J \equiv s^{-7} \pmod n = 1818, \qquad \text{check: } s^{7} \cdot J \bmod n = 1.
\]

\textbf{Round:} the prover picks $r = 37$ and sends the commitment $x = 37^{7} \bmod 3233 = 1700$. The verifier picks challenge $c = 3$. The prover, knowing $s=100$, replies with $y = 37 \cdot 100^{3} \bmod 3233 = 1548$.

The verifier checks: $y^{7} \cdot J^{3} \bmod 3233 = 1548^7 \cdot 1818^3 \bmod 3233 = 1700 = x$. \textbf{Match} — the verifier accepts, and at no point did $s=100$ itself appear in any message exchanged.
\end{example}

\begin{example}[Why a cheating prover cannot do better than guessing]
Suppose an impostor does \emph{not} know $s$, but tries to cheat by anticipating the verifier's challenge. Working backwards, they pick an arbitrary response $y=222$ and a guessed challenge $c=3$, then compute the commitment that would make the check pass \emph{for that guess}:
\[
x = y^{7} \cdot J^{3} \bmod 3233 = 937.
\]
If the verifier happens to challenge with $c=3$ (matching the guess), the check passes: $y^7 J^3 \bmod 3233 = 937 = x$. But if the verifier instead challenges with, say, $c=5$, the same fixed $(x,y)=(937,222)$ fails: $y^7 J^5 \bmod 3233 = 789 \ne 937$. Since $x$ was committed \emph{before} the challenge was revealed, the impostor had to guess $c$ in advance — correctly with probability exactly $1/v = 1/7$, matching Proposition~\ref{prop:gq-soundness}'s soundness claim precisely.
\end{example}

\begin{remark}[The bridge back to earlier material]
Guillou--Quisquater is structurally the RSA-flavored cousin of the Fiat--Shamir identification scheme (which uses $v=2$ and a modulus whose factorization is likewise the trapdoor) — the underlying commit/challenge/response skeleton (Algorithm~\ref{alg:gq} above) is the same three-move structure used by essentially every modern zero-knowledge identification and signature scheme, whatever the underlying hard problem (factoring/RSA here; discrete log for Schnorr/Fiat--Shamir; lattice problems for schemes built the same way from LWE/SIS). Recognizing this shared skeleton — commit, challenge, response, verify — is the single most transferable idea to carry out of this lecture.
\end{remark}

\newpage
\section{Summary}

\subsection*{Summary}
\begin{itemize}
\item \textbf{Padding oracles} (Bleichenbacher): a single bit of decryption-validity feedback, queried adaptively, fully decrypts PKCS\#1v1.5 ciphertexts — the flaw is in error-handling, not in RSA's algebra.
\item \textbf{Timing attacks} (Kocher): secret exponent bits leak through execution time of modular exponentiation; blinding (randomizing the input) plus constant-time arithmetic are both required as defense.
\item \textbf{Fault attacks} (Bellcore): a single corrupted CRT branch, combined with one known message/signature pair, factors $n$ outright via a single gcd computation; verify-before-release closes this completely.
\item \textbf{ROCA}: a real-world case where non-uniform, undisclosed key generation collapsed the effective search space enough for Coppersmith's method (Lesson 2) to factor otherwise-safe key sizes.
\item \textbf{OAEP}: randomized, redundancy-checked padding with an IND-CCA2 security proof (in the random oracle model) that directly rules out the algebraic and oracle-based attacks from this and the previous lecture.
\item \textbf{Guillou--Quisquater}: the same RSA algebra, repurposed as a zero-knowledge identification protocol, with soundness error $1/v$ per round instead of $1/2$ — and structurally the same three-move skeleton as the lattice-based signatures earlier in the course.
\end{itemize}

\medskip\noindent\textbf{The course-level takeaway.} Across all three RSA lectures: correctness is a CRT argument (Lesson 1); security against \emph{misuse} requires avoiding a specific list of structural traps (Lesson 2); and security in \emph{deployment} requires an entirely separate discipline — constant-time arithmetic, blinding, fault detection, verified randomness in key generation, and provably secure padding (Lesson 3). None of these three layers substitutes for the others: a mathematically flawless, correctly-used RSA implementation can still be broken in milliseconds by a fault injector, and a side-channel-hardened implementation is still broken instantly if $d$ is small or $n$ is shared. Real-world RSA security is the conjunction of all three lectures, not any one of them alone.

\section*{Tutorial Exercises}
\begin{enumerate}
\item In Bleichenbacher's attack, explain concretely why the oracle needs only to distinguish ``starts with \texttt{0x00 0x02}'' rather than fully validating the entire padding string, and why this makes the attack robust to minor implementation variations across different TLS stacks.
\item Suppose a timing attack has correctly recovered the top 6 bits of a 12-bit toy private exponent $d$ used with square-and-multiply. Describe precisely what additional measurements the attacker needs to recover bit 7, and why the attack proceeds strictly most-significant-bit first.
\item Redo the Bellcore worked example with the fault instead corrupting the $q$-branch ($s_q$) rather than the $p$-branch, and show that the same gcd technique now recovers $p$ instead of $q$.
\item Using $M = 3 \cdot 5 \cdot 7 \cdot 11 \cdot 13 = 15015$, compute the multiplicative order of $65537 \bmod M$, and compare the resulting candidate-residue count against $M$ itself and against $\sqrt{n}$ for a realistic 2048-bit RSA modulus. What does this comparison tell you about why ROCA is practical while generic factoring is not?
\item Explain why OAEP's random oracle security proof does not, by itself, guarantee security against the timing and fault attacks earlier in this lecture. What does this imply about the completeness of a security argument that only addresses the mathematical padding layer?
\item Compare Guillou--Quisquater's soundness argument to Fiat--Shamir's. If a Fiat--Shamir-style scheme uses $v=2$ per round and needs $t$ rounds to reach soundness error $2^{-t}$, how many G-Q rounds are needed to reach the same soundness error using $v = 2^{16}+1 = 65537$? What communication-cost trade-off does this reveal?
\end{enumerate}


\end{document}
